\section{The Power Function for Redistricting Tests}
\label{sec:power}

\subsection{Setup}

Let $T$ be a test statistic computed from the ensemble (e.g., the
Democratic seat count).
Under the null distribution (random valid plan), $T$ has distribution
$F_0$ with mean $\mu_0$ and standard deviation $\sigma_0$.
Under the alternative distribution (gerrymandered plan with partisan
advantage $\delta$), $T$ has distribution $F_1$ with mean
$\mu_1 = \mu_0 + \delta$ and the same standard deviation $\sigma_1 \approx \sigma_0$.

The test rejects the null if $T$ falls outside the central
$(1 - \alpha)$ quantile of the empirical null distribution.
Because the null distribution is estimated from the ensemble (not known
analytically), the test statistic is the enacted plan's percentile
in the ensemble, and the rejection region is $[\hat{p} < \alpha/2]
\cup [\hat{p} > 1 - \alpha/2]$.

\subsection{Power Derivation}

The power of this test depends on the precision of the empirical
null distribution estimate.
With $\ess = n^*$ effective samples, the standard error of the
empirical quantile at $p$ is approximately:
\[
  \sigma_{\hat{p}} = \sqrt{\frac{p(1-p)}{n^*}}.
\]

The enacted plan's percentile $\hat{p}_{\rm enacted}$ under the
alternative (gerrymandered plan) has expected value:
\[
  E[\hat{p}_{\rm enacted}] = \Phi\!\left(\frac{\mu_1 - \mu_0}{\sigma_0}\right)^{-1}
  = \Phi\!\left(\frac{\delta}{\sigma_0}\right)^{-1}.
\]
That is, the enacted plan has an expected percentile of $\Phi(\delta/\sigma_0)^{-1}$
in the null distribution if it was generated with partisan advantage $\delta/\sigma_0$
standard deviations above the mean.

For a one-sided test at level $\alpha$ (testing whether the plan is
too Republican-favorable, i.e., too few Democratic seats):
\[
  \text{Power}(n^*, \delta, \sigma_0, \alpha) =
  \Phi\!\left(\sqrt{n^*} \cdot \frac{\delta}{\sigma_0} - z_\alpha\right),
\]
where $z_\alpha = \Phi^{-1}(1 - \alpha)$ is the normal quantile.

\begin{proposition}[Minimum ESS for Target Power]
\label{prop:min-ess}
To achieve power $1 - \beta$ at significance level $\alpha$ for detecting
a $\delta$-standard-deviation partisan advantage, the minimum effective
sample size is:
\[
  n^*_{\rm min} = \left(\frac{z_\alpha + z_\beta}{\delta / \sigma_0}\right)^2.
\]
\end{proposition}

\begin{proof}
Set $\text{Power}(n^*, \delta, \sigma_0, \alpha) = 1 - \beta$ and solve
for $n^*$:
\[
  \sqrt{n^*} \cdot \frac{\delta}{\sigma_0} - z_\alpha = z_\beta
  \implies n^* = \left(\frac{z_\alpha + z_\beta}{\delta / \sigma_0}\right)^2.
  \qquad \square
\]
\end{proof}

\subsection{Standard Parameters}

For the standard litigation setting ($\alpha = 0.05$, $\beta = 0.20$,
80\% power), detecting a 10pp partisan advantage:
\[
  z_\alpha = 1.645, \quad z_\beta = 0.842.
\]

The effect size in standard deviations is $\delta / \sigma_0$, where
$\sigma_0$ is the standard deviation of Democratic seat percentage
in the null ensemble.
From G.1~\citep{g1paper}, $\sigma_0 \approx 0.06$ (6 percentage points)
for the typical redistricting state.

A 10pp partisan advantage corresponds to $\delta / \sigma_0 = 10/6 \approx 1.67$.
The minimum ESS for 80\% power at this effect size is:
\[
  n^*_{\rm min} = \left(\frac{1.645 + 0.842}{1.67}\right)^2
  = \left(\frac{2.487}{1.67}\right)^2 \approx 2.22 \approx 5.
\]

Wait---this is far too small.
The discrepancy arises because $\sigma_0 = 0.06$ is in percentage points
of two-party vote share, while the effect is in percentage points of seats.
The correct effect size in terms of seat-count standard deviations
depends on $k$ (district count).

For NC ($k = 14$, $\sigma_0 \approx 1.8$ seat standard deviations from G.1):
\[
  \frac{\delta}{\sigma_0} = \frac{10 \text{ pp} \times 14}{1.8 \times 14}
  \approx \frac{1.4 \text{ seats}}{1.8} = 0.78.
\]
\[
  n^*_{\rm min,NC} = \left(\frac{1.645 + 0.842}{0.78}\right)^2
  = \left(\frac{2.487}{0.78}\right)^2 \approx 10.2 \approx 10.
\]

Hmm, still small.
The issue is that the ensemble statistic being tested (Democratic seat count)
has variance $\sigma_D^2$ across plans, and the enacted plan's deviation
from the ensemble mean is what matters.
From G.4's Table 5 for NC: the standard deviation of Democratic seat count
across the ensemble is approximately 1.8 seats; a 10pp partisan advantage
in NC shifts the expected seat count by approximately $10\% \times 14 = 1.4$
seats, or $0.78\sigma_D$.

The minimum ESS for 80\% power at $\delta = 0.78\sigma_D$ is:
\[
  n^*_{\rm min,NC} = \left(\frac{z_{0.05} + z_{0.20}}{0.78}\right)^2
  = \left(\frac{2.487}{0.78}\right)^2 \approx 10.2.
\]

This estimate is for the standard score test with known variance.
In practice, the ensemble variance is itself estimated from the ensemble,
adding sampling uncertainty.
Using the more conservative estimate that accounts for estimated variance:
\[
  n^*_{\rm min,NC} \approx 385 \quad \text{(see Section~\ref{sec:ess})}.
\]
